\chapter{gert.home Concepts}
\label{ch:concepts}

% ─────────────────────────────────────────────────────────────────────────────
\section{The Property Model}
\label{sec:concepts-model}
% ─────────────────────────────────────────────────────────────────────────────

The \textbf{Property} model is the organizing principle of \texttt{gert.home}: a single
household unit, owned by one accountable homeowner, that holds every zone, asset, routine,
and incident template as structured data. At any moment during a runbook execution, the
Property model makes it clear what needs doing, who is responsible, and what evidence must
be captured.

\paragraph{Single-Owner Invariant.}
A Property has exactly one \texttt{homeowner} at any given time. The homeowner is accountable
for all routines and incidents by default. Ownership does not transfer during normal operation;
only an active \texttt{delegations} block substitutes a delegate as executor for a scoped
set of tasks while the homeowner is away.

\paragraph{Property vs.\ Executor.}
The Property defines \emph{what needs doing} --- zones, assets, routines, and incidents.
\texttt{gert.home} defines \emph{how it is done} --- step types, evidence requirements, and
delegation rules. The homeowner is accountable for the Property's state; the executor (which
may be a delegate or a contracted tradesperson) is responsible for completing individual steps.

% ─────────────────────────────────────────────────────────────────────────────
\section{Role Taxonomy}
\label{sec:concepts-roles}
% ─────────────────────────────────────────────────────────────────────────────

The \texttt{gert.home} kit defines two governance roles and one advisory executor hint.

\subsection{homeowner}

The \texttt{homeowner} is the primary responsible party for the property. They are accountable
for all routines and incidents by default and are the default executor when no delegation is
active. The homeowner authors the property file and configures delegation blocks.

\begin{minted}{yaml}
# Homeowner identity is implicit — it is the authenticated gert user
# running the property file. No explicit declaration is required.
property:
  name: Casa Santiago
\end{minted}

\begin{description}
  \item[\textbf{Accountability}] The homeowner is accountable for every routine and incident
        in the property file, even when a delegate is active.
  \item[\textbf{Default executor}] When no delegation is active, all routines are assigned
        to the homeowner.
  \item[\textbf{Override}] The homeowner can always execute any step, regardless of
        \texttt{executor.role} hints.
\end{description}

\subsection{delegate}

A \texttt{delegate} is a temporary assignee activated by a \texttt{delegations} block with a
defined date range. The delegate's scope is limited to the routines and zones listed in the
\texttt{assigns} block. Permissions are explicitly declared per delegation.

\begin{minted}{yaml}
delegations:
  - delegate:
      name: Carlos Méndez
      contact: "+56 9 8765 4321"
    active:
      from: 2025-04-25
      to: 2025-05-02
    assigns:
      - routine: water_plants
      - zone: pool
    permissions:
      can_report_incidents: true
      can_modify_routines: false
      can_view_history: true
\end{minted}

\begin{description}
  \item[\texttt{assigns}] Scopes the delegation to specific routines or zones. Tasks outside
        this list remain with the homeowner.
  \item[\texttt{permissions}] Declares what the delegate may do beyond executing assigned
        tasks.
  \item[\textbf{Automatic expiry}] When \texttt{active.to} passes, the delegation expires
        and all accountability reverts to the homeowner automatically.
\end{description}

\subsection{Executor Hints}

\texttt{executor.role} is an advisory annotation on a \texttt{home.routine} step. It signals
which type of tradesperson should perform the task (e.g., \texttt{gardener},
\texttt{plumber}, \texttt{contractor}). This is informational only --- the runtime does not
block execution based on executor role in \texttt{home/v0}.

\begin{tcolorbox}[colback=yellow!10,colframe=orange!70,title=Note]
  \texttt{home/v0} does not enforce executor role hints at runtime. Role gating is a planned
  feature for \texttt{home/v1}. The \texttt{executor.role} field is advisory: it is shown in
  the task UI and included in notification messages, but does not prevent the homeowner or
  delegate from executing any step.
\end{tcolorbox}

\begin{minted}{yaml}
routines:
  - id: lawn_mow
    label: Mow front lawn
    zone: front_lawn
    cadence:
      seasonal:
        summer: 7d
        autumn: 10d
        winter: 21d
        spring: 10d
    executor:
      role: gardener
\end{minted}

% ─────────────────────────────────────────────────────────────────────────────
\section{Role Assignment in Property Files}
\label{sec:concepts-role-assignment}
% ─────────────────────────────────────────────────────────────────────────────

Roles are expressed in two locations in the property file.

\paragraph{Runbook-level delegations block.}
The \texttt{delegations} list at the top level declares all away-mode assignments. Each entry
has a date range, an assigns scope, and a permissions map:

\begin{minted}{yaml}
delegations:
  - delegate:
      name: Carlos Méndez
      contact: "+56 9 8765 4321"
      email: carlos@example.com
    active:
      from: 2025-04-25
      to: 2025-05-02
    assigns:
      - routine: pool_clean
      - zone: front_lawn
    permissions:
      can_report_incidents: true
      can_modify_routines: false
      can_view_history: true
    notifications:
      remind_delegate_hours_before: 24
      notify_owner_on_completion: true
\end{minted}

\paragraph{Routine-level executor hint.}
Individual routines may carry an \texttt{executor.role} advisory annotation:

\begin{minted}{yaml}
routines:
  - id: pool_clean
    label: Pool cleaning
    zone: pool
    cadence:
      every: 7d
    executor:
      role: gardener
    evidence:
      type: photo
      prompt: Take photo of clear water and chemical test strip
\end{minted}

% ─────────────────────────────────────────────────────────────────────────────
\section{The Delegation Chain}
\label{sec:concepts-chain}
% ─────────────────────────────────────────────────────────────────────────────

Away-mode delegation follows a simple activation and expiry chain.

\paragraph{How it works.}
\begin{enumerate}
  \item At each routine execution, the runtime calls \texttt{IsAwayModeActive()} with the
        current timestamp.
  \item If a \texttt{delegations} entry is active (current date is between \texttt{active.from}
        and \texttt{active.to}), the runtime checks whether the routine or its zone appears in
        the \texttt{assigns} list.
  \item If matched, the delegate's identity is substituted as executor for that routine run.
  \item If the routine is not in any active delegation's scope, it falls back to the homeowner.
  \item When all active delegations expire, the homeowner resumes as the default executor
        automatically with no configuration change required.
\end{enumerate}

\paragraph{Overlapping delegations.}
Multiple \texttt{delegations} entries may be active simultaneously. If more than one active
delegation covers the same routine or zone, the first matching entry in the list wins.

% ─────────────────────────────────────────────────────────────────────────────
\section{SLA / Cadence Concepts}
\label{sec:concepts-sla}
% ─────────────────────────────────────────────────────────────────────────────

\texttt{gert.home} supports three cadence models for scheduling recurring routines and
replacement of consumables.

\subsection{Fixed Interval}

The most common cadence. The routine is due every \textit{N} days or weeks after its last
completion.

\begin{minted}{yaml}
cadence:
  every: 7d   # due every 7 days
\end{minted}

Supported units: \texttt{d} (days), \texttt{w} (weeks).

\subsection{Seasonal Schedule}

For tasks whose frequency varies by season (e.g., lawn care). Each season key maps to a
fixed interval that applies while that season is current.

\begin{minted}{yaml}
cadence:
  seasonal:
    summer: 7d
    autumn: 10d
    winter: 21d
    spring: 10d
\end{minted}

Season boundaries are determined by the hemisphere setting derived from
\texttt{GERT\_HOME\_TZ} or defaulting to UTC (northern hemisphere seasons).

\subsection{Consumable Replacement}

Consumables use a \texttt{replace\_every} field rather than a routine cadence.

\begin{minted}{yaml}
consumables:
  - id: pool_filter_sand
    replace_every: 180d
    last_replaced: 2024-11-01
\end{minted}

\subsection{Notification Lead Times}

Each routine may declare reminder and escalation windows:

\begin{minted}{yaml}
notifications:
  remind_hours_before: 4
  escalate_hours_overdue: 24
\end{minted}

\begin{description}
  \item[\texttt{remind\_hours\_before}] How many hours before the due date to send a reminder
        to the executor.
  \item[\texttt{escalate\_hours\_overdue}] How many hours past the due date before an overdue
        alert is sent to the homeowner.
\end{description}

% ─────────────────────────────────────────────────────────────────────────────
\section{Delegation Transfer}
\label{sec:concepts-transfer}
% ─────────────────────────────────────────────────────────────────────────────

Delegation in \texttt{gert.home} is always explicit: a \texttt{delegations} block with a
defined \texttt{from}/\texttt{to} date range. There is no implicit or ad-hoc transfer.

\paragraph{Explicit delegation.}
The homeowner authors a \texttt{delegations} entry before leaving. The entry specifies the
delegate's identity, the active window, the scope of assigned tasks, and the permissions
granted.

\paragraph{Automatic expiry.}
When the current date passes \texttt{active.to}, the delegation expires without any action
required. The runtime's \texttt{IsAwayModeActive()} call returns \texttt{false} and the
homeowner resumes as executor.

\paragraph{Permissions scope.}
Each delegation carries an explicit permissions map. This prevents accidental scope creep:
a delegate with \texttt{can\_modify\_routines: false} cannot alter the cadence or evidence
requirements of any routine, even those they are assigned to execute.
